\chapter{A timeline of DNA computing}\label{app:timeline}
\markboth{A timeline of DNA computing}{A timeline of DNA computing}

The milestones below trace how DNA computing grew from a single experiment into
a family of research programmes. Figure~\ref{fig:timeline} groups them into five
themes; the table gives one primary reference for each. The selection is
necessarily partial and favours results that the chapters of this book build
on. A milestone records what was demonstrated, under the conditions reported,
not a general capability.

\TikzPlate{backmatter/figures/timeline}{Milestones of DNA computing, grouped
into five themes. Read down a lane to follow one research thread; read across a
row to see what happened in the same year. Early work searched for solutions
by filtering molecular libraries; since about 2006 the field has increasingly
programmed molecular \emph{dynamics} (strand displacement, self-assembly) and,
more recently, molecular \emph{learning} and information storage.}{fig:timeline}

\begin{center}
\small
\begin{longtable}{@{}p{0.08\linewidth}p{0.62\linewidth}p{0.22\linewidth}@{}}
\toprule
Year & Milestone & Reference\\
\midrule
\endhead
1953 & Double-helical structure of DNA proposed & \cite{WatsonCrick1953}\\
1961 & Logically irreversible operations must dissipate heat & \cite{Landauer1961}\\
1982 & Reversible computation; polymerase as a Brownian computer & \cite{Bennett1982}\\
1982 & Branched DNA junctions proposed as building blocks for lattices & \cite{Seeman1982}\\
1988 & PCR with a heat-stable polymerase & \cite{Saiki1988}\\
1994 & Directed Hamiltonian path solved with DNA & \cite{Adleman1994}\\
1995 & Satisfiability by molecular filtering & \cite{Lipton1995}\\
1996 & Computational power of DNA filtering models analysed & \cite{Boneh1996}\\
1997 & Maximal clique problem solved with DNA & \cite{Ouyang1997}\\
1998 & Sticker model of DNA memory & \cite{Roweis1998}\\
1998 & Two-dimensional DNA tile crystals self-assemble & \cite{Winfree1998}\\
1998 & Unified nearest-neighbour thermodynamic parameters & \cite{SantaLucia1998}\\
2000 & DNA-fuelled molecular tweezers & \cite{Yurke2000}\\
2000 & Synthetic genetic toggle switch and oscillator in bacteria & \cite{Gardner2000,ElowitzLeibler2000}\\
2001 & Autonomous enzyme-driven finite automaton & \cite{Benenson2001}\\
2002 & 20-variable 3-SAT instance solved with DNA & \cite{Braich2002}\\
2003 & Deoxyribozyme automaton plays tic-tac-toe & \cite{Stojanovic2003}\\
2004 & Algorithmic self-assembly of Sierpinski triangles & \cite{Rothemund2004}\\
2004 & Molecular automaton logically controls gene expression in vitro & \cite{Benenson2004}\\
2006 & DNA origami: a scaffold folded by staple strands & \cite{Rothemund2006}\\
2006 & Enzyme-free nucleic-acid logic circuits & \cite{Seelig2006}\\
2007 & Entropy-driven catalytic strand-displacement networks & \cite{Zhang2007}\\
2009 & Toehold exchange controls displacement kinetics & \cite{ZhangWinfree2009}\\
2010 & DNA can implement arbitrary chemical reaction networks & \cite{Soloveichik2010}\\
2011 & Seesaw circuits compute square roots of 4-bit numbers & \cite{QianWinfree2011}\\
2011 & Hopfield-type associative memory from DNA & \cite{QianWinfreeBruck2011}\\
2012 & Megabit-scale digital data written into DNA & \cite{Church2012}\\
2012 & Logic-gated DNA nanorobot; DNA-brick 3D structures & \cite{Douglas2012,Ke2012}\\
2013 & Archival DNA storage with error-tolerant encoding & \cite{Goldman2013}\\
2013 & Recombinase logic and memory in living cells & \cite{Siuti2013,Bonnet2013}\\
2013 & Programmable chemical controllers made from DNA & \cite{Chen2013}\\
2016 & Design automation for genetic circuits & \cite{Nielsen2016}\\
2017 & Fountain-code DNA storage; a movie recorded in bacterial genomes & \cite{ErlichZielinski2017,Shipman2017}\\
2017 & Enzyme-free nucleic-acid oscillator; cargo-sorting DNA robot & \cite{Srinivas2017,Thubagere2017}\\
2017 & Ribocomputing devices perform logic in cells & \cite{Green2017}\\
2018 & Winner-take-all DNA neural networks classify handwritten digits & \cite{CherryQian2018}\\
2018 & Random access in large-scale DNA storage & \cite{Organick2018}\\
2018 & Design principles for leakless strand displacement & \cite{Wang2018}\\
2019 & One reprogrammable tile set runs 21 different algorithms & \cite{Woods2019}\\
2020 & DNA molecular computation classifies clinical samples & \cite{Zhang2020cancer}\\
2022 & Molecular convolutional neural networks & \cite{Xiong2022}\\
2022 & DNA Typewriter: sequential recording in human cells & \cite{Choi2022}\\
2023 & DNA-based programmable gate arrays & \cite{Lv2023}\\
2024 & Pattern recognition by competitive nucleation of tiles & \cite{Evans2024}\\
2024 & Parallel storage with epigenetic bits & \cite{ZhangC2024}\\
2025 & Supervised learning carried out by DNA in vitro & \cite{CherryQian2025}\\
2025 & Heat-rechargeable DNA circuits and neural networks & \cite{SongQian2025}\\
2026 & Thermodynamically favoured scaffolded DNA computer & \cite{Sterin2026}\\
\bottomrule
\end{longtable}
\end{center}
